\section{Discussion}
\label{sec:discussion}

\subsection{Interpretation}
In this corpus and for this detector, the data give no evidence that LLM paraphrasing lowers recall, and they bound the possible loss at about 3\% of e-mails (Section~\ref{sec:results}). The score did fall by 0.53 on average, so the paraphraser changes what the detector sees. Recall stayed unchanged because the median original e-mail lies 4.27 score units above the threshold, while a shift that costs 2 percentage points is about $-1.10$ (Eq.~\eqref{eq:breakeven}); the observed shift is roughly half of it. This suggests a large safety margin on these data, not that paraphrasing is harmless in general. The model of Eq.~\eqref{eq:shift} slightly overestimates recall (error 1.2 percentage points); normality is an approximation, and it would be weaker for saturated scores such as neural probabilities.

The margin is probably so large because of a ceiling effect. ROC-AUC is about 0.9997, and the legitimate e-mail (2008 mailing lists and personal mail) differs strongly in style from the phishing e-mail (modern marketing and wallet-themed HTML mail). The detector may rely partly on corpus style rather than on phishing semantics. Two label giveaways were removed before splitting and the result did not change, but residual style differences remain, and a small paraphrasing effect could be hidden by the ceiling. Thapa et~al.\ report that LLM-rephrased e-mails degrade detector performance~\cite{thapa_phishing_2025}. Our result for a TF-IDF and logistic-regression detector does not reproduce a drop: recall is unchanged and the possible loss is bounded at 3\%. We do not compare numbers, because we have only read their summary and their detectors, data and metric may differ from ours; a likely reading is that a classical detector with a wide margin on easy data is less sensitive than the detectors and data they tested, which our data cannot confirm.

\subsection{Limitations}
\emph{Internal validity.} There are 97 pairs of short e-mails (160 to 900 characters), so the Wilson interval of recall is wide, [92.8\%, 99.4\%]. The quality filter checks length and link retention, not meaning; no independent human check of the paraphrases was made. \emph{External validity.} Only one detector (TF-IDF and logistic regression) and one paraphraser (Claude; the planned Qwen2.5-1.5B-Instruct was unavailable) were used. The data are English, one phishing source and one legitimate source of a different era. DistilBERT, the zero-shot LLM, the regression of RQ2 and the augmentation of RQ3 were not run, so nothing is claimed about them. \emph{Construct validity.} Recall at FPR~1\% is a single operating point, and only the message text is scored, not headers, authentication or attachments.

\section{Conclusion}
\label{sec:conclusion}
We measured how much LLM paraphrasing lowers the recall of a text-based phishing detector at a fixed false-positive rate of 1\%. On 97 paired e-mails, recall was 0.979 before and after paraphrasing, no e-mail was lost (upper 95\% bound 3.0\%), and a location-shift model explained why: the score fell by 0.53 but the median margin of 4.27 absorbed it. The finding is limited to one detector, one paraphraser, short e-mails and an easily separable corpus pair, and it does not show that paraphrasing is harmless for detectors trained on harder data. Future work is to run the neural detectors on a GPU, to rebuild the corpus with modern same-source legitimate e-mail, and to paraphrase the training set for RQ2 and RQ3.
